\documentclass{article}

\usepackage{amsmath,amssymb}
\usepackage{xurl}
\usepackage[hidelinks]{hyperref}

% Shared commands for notes in docs/paper-gaps/.

\DeclareUnicodeCharacter{2080}{\ifmmode{}_{0}\else\textsubscript{0}\fi}
\DeclareUnicodeCharacter{2081}{\ifmmode{}_{1}\else\textsubscript{1}\fi}
\DeclareUnicodeCharacter{2082}{\ifmmode{}_{2}\else\textsubscript{2}\fi}
\DeclareUnicodeCharacter{2099}{\ifmmode{}_{n}\else\textsubscript{n}\fi}

\newcommand{\Lean}{\textsc{Lean}}
\ExplSyntaxOn
\NewDocumentCommand{\leanid}{m}{
  \ifmmode
    \textsf{\detokenize{#1}}
  \else
    \group_begin:
    \urlstyle{sf}
    \tl_set:Nn \l_tmpa_tl {#1}
    \tl_replace_all:Nnn \l_tmpa_tl {\_} {_}
    \exp_args:NV \path \l_tmpa_tl
    \group_end:
  \fi
}
\ExplSyntaxOff
\providecommand{\tr}{\operatorname{tr}}
\newcommand{\tnleanIssueBase}{https://github.com/LionSR/TNLean/issues/}
\newcommand{\tnleanPrBase}{https://github.com/LionSR/TNLean/pull/}
\newcommand{\tnleanDocBase}{https://sirui-lu.com/TNLean/docs/}
\newcommand{\ghissue}[1]{\href{\tnleanIssueBase#1}{GitHub issue \##1}}
\newcommand{\ghpr}[1]{\href{\tnleanPrBase#1}{PR \##1}}
\newcommand{\doclink}[1]{\href{\tnleanDocBase#1.html}{\path{#1}}}

% Dirac notation and the identity operator, as in blueprint/src/macros/common.tex.
% \providecommand keeps note-local definitions authoritative.
\usepackage{dsfont}
\providecommand{\ket}[1]{|#1\rangle}
\providecommand{\bra}[1]{\langle #1|}
\providecommand{\braket}[2]{\langle #1 | #2 \rangle}
\providecommand{\ketbra}[2]{|#1\rangle\!\langle #2|}
\providecommand{\Id}{\mathds{1}}
\providecommand{\one}{\mathds{1}}
\providecommand{\C}{\mathbb{C}}
\providecommand{\MN}[1]{M_{#1}(\C)}
\providecommand{\MD}{M_D(\C)}

% External referencing for paper sources.  The build helper writes sanitized
% auxiliary files under build/paper-aux/ and excludes duplicated source labels.
\usepackage{xr}

% Import a generated paper auxiliary file with a caller-supplied label prefix.
% The candidate paths support builds from docs/paper-gaps/, the repository root,
% and build/paper-aux/ itself.
\newcommand{\importpaperaux}[2]{%
  \IfFileExists{../../build/paper-aux/#2.aux}{%
    \externaldocument[#1]{../../build/paper-aux/#2}%
  }{%
    \IfFileExists{build/paper-aux/#2.aux}{%
      \externaldocument[#1]{build/paper-aux/#2}%
    }{%
      \IfFileExists{#2.aux}{%
        \externaldocument[#1]{#2}%
      }{}%
    }%
  }%
}

% General external reference macros
\newcommand{\paperref}[2]{\ref{#1-#2}}
\newcommand{\papereqref}[2]{(\ref{#1-#2})}

% CPSV16 source (1606.00608 / MPDO-22-12-17-2.tex) external document and shortcuts
\importpaperaux{cpsv16-}{1606.00608/MPDO-22-12-17-2-tnlean}
\newcommand{\cpsvref}[1]{\paperref{cpsv16}{#1}}
\newcommand{\cpsveqref}[1]{\papereqref{cpsv16}{#1}}

% Verdict marker: \gapnote{<kind>}{<status>}. Kind is the policy
% classification (clarification, local-correction, scope-restriction,
% unfaithful, false-source, open-gap); status is open, wip (elimination
% underway), resolved, or historical. Machine-read by the site index;
% prints a verdict line.
\newcommand{\gapnote}[2]{\par\noindent\textbf{Verdict:} #1 (#2).\par}


\title{The Periodic-Sector Argument in Proposition 4.13}
\author{}
\date{2026-07-13}

\begin{document}
\maketitle
\gapnote{local-correction}{resolved}

\section{The printed argument}

The proof of Proposition~4.13 of
Ref.~\cite[Proposition~4.13]{Cirac2016MPDO_arXiv} excludes nontrivial periodic
sectors of the vertically viewed tensor. The relevant source passage is
\path{Papers/1606.00608/MPDO-22-12-17-2.tex}, lines~1888--1893. It asserts
that a nontrivial period provides an orthogonal projector $Q$ such that, for
every $N$,
\begin{align}
    Q_1H^{(N+1)}
        &\ne H^{(N+1)}Q_1,
    \label{eq:cpgsv17_periodic_sector_projector_printed_noncommutation}\\
    Q_1[H^{(N+1)}]^p
        &=[H^{(N+1)}]^pQ_1,
    \label{eq:cpgsv17_periodic_sector_projector_printed_power_commutation}
\end{align}
where
\begin{align}
    Q_1
        &=Q\otimes\Id^{\otimes N}.
    \label{eq:cpgsv17_periodic_sector_projector_first_site_projector}
\end{align}
Because $H^{(N+1)}$ is positive semidefinite,
Eq.~\ref{eq:cpgsv17_periodic_sector_projector_printed_power_commutation}
implies
\begin{align}
    Q_1H^{(N+1)}
        &=H^{(N+1)}Q_1,
    \label{eq:cpgsv17_periodic_sector_projector_power_implies_commutation}
\end{align}
contradicting
Eq.~\ref{eq:cpgsv17_periodic_sector_projector_printed_noncommutation}.

\section{The quantifier correction}

Let $\widetilde M$ denote the tensor viewed in the vertical direction. The
cyclic decomposition of its peripheral spectrum gives an orthogonal projector
$Q$ that is invariant under every word of one full period but is displaced by
one vertical layer:
\begin{align}
    Q\widetilde M
        &\ne Q\widetilde M Q.
    \label{eq:cpgsv17_periodic_sector_projector_vertical_displacement}
\end{align}
The representative-grouped form of Lemma~L in
Ref.~\cite{Cirac2016MPDO_arXiv} proves the contrapositive implication
\begin{align}
    Q\widetilde M\ne Q\widetilde M Q
        &\Longrightarrow
        \exists N\geq0:\
        Q_1H^{(N+1)}\ne H^{(N+1)}Q_1.
    \label{eq:cpgsv17_periodic_sector_projector_exists_noncommuting_length}
\end{align}
It does not prove noncommutation for every $N$, as printed at source
line~1889. The stronger quantifier is unnecessary. The positivity argument at
source lines~1890--1893 is pointwise in $N$, so the witness supplied by
Eq.~\ref{eq:cpgsv17_periodic_sector_projector_exists_noncommuting_length}
suffices.

This change is a local correction to the intermediate statement. The corrected
periodic-sector implication is proved both from literal CPSV canonical form and
from the stronger normalized BNT-refined horizontal form. It completes the
periodic-exclusion step in Proposition~4.13 of
Ref.~\cite{Cirac2016MPDO_arXiv}. The subsequent vertical phase-class grouping,
gauge normalization, and coisometry construction are recorded in
\path{cpgsv17_vertical_cf_grouping.tex}.

\section{Formal proof}

The literal-CPSV proof of
Eq.~\ref{eq:cpgsv17_periodic_sector_projector_exists_noncommuting_length} is
\leanid{MPSTensor.IsCPSVCanonicalForm.exists_not_commute_of_displaced} in
\doclink{TNLean/MPS/MPDO/CPSVPeriodicExclusion}. Suppose, for contradiction,
that $Q_1$ commutes with $H^{(N+1)}$ for every $N$. Idempotence of $Q$ then
gives
\begin{align}
    Q_1H^{(N+1)}
        &=Q_1H^{(N+1)}Q_1.
    \label{eq:cpgsv17_periodic_sector_projector_sandwiched_commutation}
\end{align}
Write both sides of
Eq.~\ref{eq:cpgsv17_periodic_sector_projector_sandwiched_commutation} as
first-site contractions of the doubled-index tensor. The original-space form
of Lemma~L then yields
\begin{align}
    Q\widetilde M
        &=Q\widetilde M Q,
    \label{eq:cpgsv17_periodic_sector_projector_lemma_l_invariance}
\end{align}
contradicting
Eq.~\ref{eq:cpgsv17_periodic_sector_projector_vertical_displacement}. The
corresponding theorem under the stronger normalized BNT-refined horizontal
hypothesis is
\leanid{MPOTensor.IsHorizontalCF.exists_not_commute_of_displaced} in
\doclink{TNLean/MPS/MPDO/HorizontalBNT}.

The cyclic decomposition and its letter-shift property are formalized in
\path{TNLean/MPS/MPDO/CyclicProjector.lean}. They provide a period $p\geq1$,
an orthogonal projector $Q$, invariance under full-period words, and
displacement by one letter. Choose the length $N+1$ supplied by
Eq.~\ref{eq:cpgsv17_periodic_sector_projector_exists_noncommuting_length}.
Full-period invariance transfers to the $p$-fold vertically stacked tensor and
gives
\begin{align}
    Q_1[H^{(N+1)}]^p
        &=[H^{(N+1)}]^pQ_1.
    \label{eq:cpgsv17_periodic_sector_projector_witness_power_commutation}
\end{align}
The positive-semidefinite power-commutation theorem then gives
\begin{align}
    Q_1H^{(N+1)}
        &=H^{(N+1)}Q_1,
    \label{eq:cpgsv17_periodic_sector_projector_witness_commutation}
\end{align}
which contradicts the choice of $N$.

The literal conclusion is
\leanid{MPOTensor.hasNoPeriodicVectors_verticalTensor_of_cpsvCanonicalForm} in
\doclink{TNLean/MPS/MPDO/CPSVPeriodicExclusion}. The stronger-hypothesis
variant is
\leanid{MPOTensor.hasNoPeriodicVectors_verticalTensor_of_horizontalCF} in
\doclink{TNLean/MPS/MPDO/CyclicProjector}. Neither theorem adds an assumption
beyond its stated canonical-form and MPDO hypotheses.

\section{Status}

The periodic-sector step is complete under literal CPSV canonical form and
under the stronger normalized BNT-refined predicate
\path{MPOTensor.IsHorizontalCF}. The all-length condition printed at source
line~1889 --- that a displaced orthogonal projector fails to commute with
$H^{(N)}$ at every length --- is used by neither proof and is not formalized.

Combining literal periodic exclusion with literal invariant-projector closure
produces normal vertical corners with positive weights and exact letterwise
reconstruction in
\leanid{MPSTensor.IsCPSVCanonicalForm.exists_normal_verticalBlockDecomp_with_isometry}.
The subsequent phase-class grouping, normalization of the grouped
representatives, and construction of the final coisometry are also complete.
Together, these results establish the literal theorem
\leanid{MPOTensor.verticalCF_of_cpsvCanonicalForm} in
\doclink{TNLean/MPS/MPDO/CPSVVerticalCanonicalForm}. The stronger theorem from
normalized BNT-refined horizontal form remains separate.

This completion concerns only Proposition~4.13 of
Ref.~\cite{Cirac2016MPDO_arXiv}. It proves neither the positive-fusion
statement in Appendix~C.4 nor the renormalization fixed-point characterization
in Theorem~4.14 of the same reference.

\paragraph{Verdict.}
The existence of one noncommuting length replaces the printed all-length
inequality. This local correction suffices for periodic-sector exclusion.
Together with the independent grouping, normalization, and coisometry results,
it completes literal Proposition~4.13.

\bibliographystyle{alpha}
\bibliography{references}

\end{document}
